\section{The Cross and the Eagle}

\textit{The Secret Language of the Fedeli D'Amore} by \textbf{Arturo Reghini} (Pietro Negri) originally appeared in \textit{Introduction to Magic}, volume 2.

\begin{quotex}
\textbf{Arturo Reghini}, \textbf{Guido De Giorgio}, and \textbf{Julius Evola} each tried to recapture the ideal of the Roman Tradition in a different way. For example, where De Giorgio saw in Dante the amalgamation of the Eagle and the Cross through a development, Reghini, as a pagan, Mason, and Pythagorean, sought to cleanly separate the two. In this essay, Reghini, as did \textbf{Rene Guenon}, counts on the work of \textbf{Luigi Valli} to explicate the hidden meaning of \textbf{Dante}`s Divine Comedy, not to mention the group of poets known as the \textbf{Fedeli d'Amore}.

Reghini reveals that the object of Dante's ``love" is really the Divine Sophia. He brings up the ideas of the active and potential intellect, a topic of great interest at that time in the Middle Ages. Astute readers will be able to relate this to the meditation formulated by \textbf{Valentin Tomberg} in the second letter, which we have mentioned several times. Reghini also points to an interesting correspondence with an idea from the The Tibetan Book of the Dead. 

\end{quotex}
Several years ago \textbf{Luigi Valli} published La Chiave della Divina Commedia [The Key to the Divine Comedy] in which, proceeding successfully along the interpretative line divined by Foscolo and then followed by \textbf{Gabriele Rossetti}, Perez, Pascoli, and a few others, he succeeded at highlighting thirty agreements between the Eagle and the Cross in the sacred poem and finds, at least in part, the doctrine hidden under the veil of the strange verses [delli versi strani]. The thought exposed and simultaneously hidden by Dante would be, very synthetically, this: The Cross showed itself impotent to redeem in fact humanity and cannot redeem it alone. The involvement of the Eagle is necessary, i.e., for authority and imperial justice, it is necessary to reestablish the Empire, take away from the Church the unfavorable donation of Constantine; the corruption of the Church and humanity will then certainly have an end, thanks to the double virtue of the Cross and the Eagle and it will actually be able to save itself. Dante proclaimed openly that on the cathedral of Saint Peter stood the unworthy usurpers, the preachers of gossip, who did not possess the genuine intention given by Christ to his first monks; and he covertly added that on the chair of the Church was seated the apocalyptic whore, he recognized the failure of the preaching of the Cross and the necessity of the intervention of the imperial Eagle to save humanity. This bold, and for certain not very orthodox from the Catholic perspective, conception inspired not only Dante's writings but also his action, understood as carrying out his program first by means of the armies of the \textbf{Templars}, and then of the Emperor.

Following logically the thread of these studies, Luigi Valli next published an extremely important, interesting, and powerful volume, entitled: Il Linguaggio segreto di Dante and dei Fedeli d'Amore [The Secret Language of Dante and the Fedeli d'Amore]. The first centuries of Italian literature and all the history and the battles of those times are the object of this study, and are presented under a light and an aspect that even now is unsuspected and unexpected. With a patient, methodic, scientific, and imposing work, Valli, resuming the misunderstood and neglected work of Rossetti, confirms and demonstrates the existence from the beginning of Italian literature of a secret language, the jargon of the Fedeli d'Amore; he deciphers its meaning, the numerous doctrinal, sectarian, and political allegories and brings back into the light a whole grandiose movement, inspired by the ``initiatic tradition" and bitter enemy of the Church of Rome.

Unable to even succinctly summarize the events of this great battle, we will only say how, through this understanding, the poets of love, the writer of the ``sweet new style", who strangely seemed to lose themselves singing of their absurd, self-conscious and inconsistent love, are transfigured into formidable battlers, into ardent champions of their Holy Faith. They dramatically tower over all the most noble figures of Cecco d'Ascoli and Dante Alighieri, who are greater, the more they are understood. We express to Luigi Valli our admiration and our recognition; his work constitutes, as we intended to point out, a ``piece of gelatin", and as much against it as the myopic and lazy misoneism and ``positive criticism", the vestal of pure aesthetic, and the shrewdness of the curious, has coalesced, the light is by now made and will end up by standing out.

The love for which the heart of Fedeli d'Amore was burning, is similar to the mystical love of Persian literature and that of the Song of Songs. Gabriel Rossetti reconnected it above all with Platonic love, which would give a pagan character to the movement. Valli demonstrates that the ``rose", the ``flower", the ``woman", (1) which is under various names the only object of this love, is the active intellect, that loves of itself the potential intellect (2): it is, as Dino Compagni sings:

\begin{table}[h]\footnotesize
\begin{tabularx}{\textwidth}{XX}
L'amorosa, Madonna Intelligenza &
The lovely Lady intelligence\\
Che fa nell'alma la sua residenza &
Who makes her home in the soul\\
Che co' la sua bielta m'ha innamorato &
Who with her beauty enthralled me
\end{tabularx}
\end{table}
To the accumulation of the proofs that Valli discovered or brought back in this regard, more of them could be added; this, for example: Dante from the principle of the Comedy speaks of the

\begin{table}[h]\footnotesize
\begin{tabularx}{\textwidth}{XX}
Divina potestate &
The divine power\\
La somma sapienza e il primo amore, &
Wisdom in the Highest and Primal Love
\end{tabularx}
\end{table}
He places his ``love" in a triad that corresponds perfectly—in the Kabbalah—to the triad of the highest sephiroth: Kether, Chokhmah, Binah, or the Crown, Wisdom, Intelligence.

If this is the woman, the domina, of the Fedeli d'Amore, it is perfectly logical that Francesco da Barberin in his Documenti di Amore puts docilitas, docility (from docere, to teach), first among the twelve virtues that Love must awaken in the novices. The tradition that puts this docility among the first requisites of initiation is transmitted down to us, as is shown for example by what Arturo Reghini writes on pages 106-108 of his book on the Parole sacre e di Passo [Sacred Words and Passage]. Even the world discipline has the double meaning of science and constraint; and the German gelebrig corresponds through its polysemy to the Latin docilis.

The transmission of the secret language of the Fedeli d'Amore in that of later sects and movements was recognized, in addition to Valli and before him, by Rossetti and Aroux, who actually pushed too much on this way and were sometimes led astray by the intent of wanting to recognize the concordance between the various sectarian jargons; but the concordance undoubtedly exists in part, and leads to posing the problem of the transmission, not of only sectarian jargon, but of the same traditional doctrine.

We, too, with Valli, believe that Rossetti, the first systematic discoverer of the sectarian jargon of the Fedeli d'Amore, was led to his interpretation by the knowledge of ancient secret traditions. If memory does not fail us, his Mystery of Platonic Love in the Middle Ages was dedicated to B. L., which is very plausibly \textbf{Bulwer Lytton}, the author of Zanoni, who, beyond having a profound esoteric erudition, was also an expert on Italian language and literature. One could perhaps think that Rossetti was inducted and initiated by Bulwer Lytton into the systematic study of the sectarian medieval jargon, a study happily taken up by Valli, who succeeded in emending, extending, and completing the results achieved by Rossetti in the last century.

We saw that Love is the ``Active Intellect"; it is, as Dante says in the last verso of the Comedy: ``The love that moves the Sun and the other stars". In the potential intellect of the Fedeli d'Amore this active intellect is awakened and operative, in the profane it is sleeping and inoperative. According to Valli, in the sectarian jargon, sleeping consistently means to be in error, to be far from the truth and in particular to belong to the Church of Rome. It is the symbolism adopted by Dante in the last cantos of the Purgatorio, in which after the immersion in the river Lethe, the river of the dream and oblivion, the immersion into Eunoe follows, by virtue of which, like a new plant (neophyte) with renewed frond. Dante becomes pure and ready to jump to the stars, that is, capable of rising to the ``kingdom of Heaven". As we noted, it is about a pagan symbolism adopted by Virgil and Plato, and that is found again in the very old Orphism and in the Eleusinian mysteries; here at the river Lethe, which sweeps away the knowledge of men, is in contrast to the fresh arising of Memory or the mnemonic virtue of the pomegranate, that gives awakening and immortality. The Platonic anamnesis, the record, is identified to the consciousness and correspondingly the truth, the aletheia, and is the negation, the passing of Lethe. The attainment of the truth is a conquest of consciousness above the dream and death; it is necessary to maintain the continuity of consciousness even through the dream and death.

Love in the initiatic sense has therefore the capacity to take away the dream and death, giving to the Fedel d'Amore a new life. That is reached through degrees of successive development.

\paragraph{The Arrow and the Rose}

\begin{quotex}
In the concluding part, Reghini points to several other Hermetic symbols found in \textbf{Dante} and other writings of the \textbf{Fedeli d'Amore}. These include the Arrow, the Rose, the Rebis, the Azoth, and the Phoenix. Finally, he mentions the significance of the analogical or spiritual meaning of such texts, which simply cannot be understood by those not prepared for it.

\end{quotex}

In \textbf{Francisco da Barberino}'s \textit{Documenti d'Amore}\footnote{\url{https://archive.org/details/DocumentiAmoreEgidiVol1/page/n18/mode/2up}} [Documents of Love], the Fedeli d'Amore is represented in the first degrees as pierced by the arrow of Love and in the last degree with some roses in the hand. The symbolism of the arrow is also found in one of the twelve figures of Basil Valentine's Azoth\footnote{\url{https://en.wikipedia.org/wiki/Basil_Valentine}\newline\url{https://www.alchemywebsite.com/prints_series_val_azoth.html}}. But the similarity between the symbolism of love and Hermetic symbolism and the link between the two traditions again turn out to be shown more through the presence of the Hermetic Rebis in one of the designs that illustrate Barberino's Documents of Love. The Rebis, or Hermetic androgyne, is a characteristic and very important Hermetic symbol, whose history we briefly treated in another work, Un codice alchemico italiano; the figure of the Rebis reproduced by Valli goes back to Dante's time and is older by several centuries than what we tracked down in the books on Hermetism.

Other concordances with the symbolism and Hermetic terminology are found in the verses of an obscure poet of love, \textbf{Nicolo dei Rossi}, who in one of his lyrics expresses ``the degrees and the virtue of true love". There are four degrees: the first is called liquefatio which is opposite, says dei Rossi, to congelazione. The second degree is called languor, the third zelus, and in the fourth, love reaches the hierarchical summit by means of ecstasy or excessus mentis. We understand therefore how one of the most important works of the literature of Love, the Romance of the Rose, (whose Italian version, Il Fiore, is due to a Florentine named \textbf{Durante} who is almost certainly \textbf{Dante}), treated alchemy explicitly and is classified in the alchemical literature. This rose sung with such moving harmony by all these poets, starting from \textbf{Ciullo l'Alcamo}, the dantesque rose candida, is clearly similar, if not identical, to the Hermetic rose of the Rosicrucians.

An important confirmation of this assimilation and affinity between Hermetism and the Fedeli d'Amore is given here by the four so-called ``Templar degrees" of Masonry which arose in France or in Germany toward the middle of the 18th century. It is about the Princes of Mercy, called also the Knights of the Sacred Delta, and also designated in another way. Their task, says the ritual, is

\begin{quotex}
to guard with fidelity the treasure of traditional wisdom, always concealing it from those who do not know how to penetrate into the third heaven.

\end{quotex}
The third heaven is the name of their temple and is, as everyone knows, the heaven of Venus. We note moreover that in Orphism and Pythagoreanism, the third heaven is the last. \textbf{Philolaus} in fact says that there are three heavens: Uranus, the Cosmos, and Olympus. The third heaven, Olympus, is the home of the gods, and Saint Paul referred to this Orphic-Pythagorean classification when he told of having been raptured to the third heaven.

Now the ``intellect" of \textbf{Dion Compagni}, writes Valli,

\begin{quotex}
stays in a palace where different locations represent probably degrees of initiation, and in that palace the third place is the salutatorio … referring us back to the frequent allusions to the third heaven or the third degree, that in the material heaven is the heaven of Venus, but in the symbols signified rather probably the worship or a higher degree of his initiation.

\end{quotex}
The Princes of Mercy by means of their triple virtue succeed in lifting up the veil of truth; and are therefore called beni emeth, the sons of Truth. Among the characteristic symbols of the degree the Palladium of the Order appears, or the statue of Truth, naked and covered with a tricolored veil. These three colors that reappear in the decorations of the Temple and in other symbols of the degree are green, white, and red, the three Hermetic colors with which Dante adorns his Beatrice (Purg, XXX, 31-33).

The numeric symbolism of the degree is based on the number three and its powers: the sacred or luminous Delta is one of its principle symbols. The word emeth, truth, contains three letters, the first, the middle, and the last of the Hebrew alphabet. Its numeric value is 441, or nine. On the throne are nine lights. In the temple are nine columns, each one of which bears a candelabra with nine lights or in all there are 81 lights. The age of 81 years is the ritual age. We will not dwell on recalling the importance Dante attached to three and nine, and with the frequency the number nine recurs in the Vita Nuova: Valli relates some verses in which \textbf{Giacomo da Lentini} proposes that ``the mercies are strict… nor by the lovers called finally who completes nine years".

As to the number 81, Valli already reported the following strange and bold passage from Dante that he writes precisely in the Convivio:

\begin{quotex}
Plato, from whom one can very well say that he had matured … living 81 years … And I believe that if Christ had not been crucified and had lived out the span which his life, according to its nature, might have passed from the mortal body into the eternal in his eighty-first year (IV, xxiv)

\end{quotex}
that is, if he had reached the ritual age of the Knights of the sacred Delta. Naturally Dante in the Vita Nuova had Beatrice die in the ninth day of the month of June in 1281; and he took care to specify that in Syria the month of June is the ninth, and that Beatrice was dead when ``the perfect number was completed nine times" in the third tenth century, or in 1281.

Among the symbols of this degree that are reconnected to the symbolism of the ``Fedeli d'Amore" we note again the arrow that was on the throne of the Most Excellent (the president of the chamber), that is obviously the arrow that Francesco da Barberino puts in the hand of Love in the first figure of his Documents of Love. This arrow is of white wood and has feathers colored partly in green and partly in red, with a gold point.

Another symbol of the degree is composed of two arrows, the two arrows of the Love of tradition, one of gold, the other of lead: the two arrows of the dantesque lyrical poem: ``Three women have come around my heart". For fuller information about this topic we refer to the Manual of \textbf{Andres Cassard}. And finally it is necessary to note how the sole Phoenix, about which there is a continual mention in the poetry of the Fedeli d'Amore and that, as Valli shows, represents the organization and the initiatic tradition always being reborn, either one or the other of the most important symbols of Hermetism, the symbol of the Rubedo. The purple Phoenix is reborn and lives among the flames of the ``philosophical fire", as the Fedelt d'Amore, burning with holy zeal (the zelus of Nicolo dei Rossi), is reborn to new life by means of the excessus mentis.

Numerous other comparisons could be established between the sectarian jargon deciphered by Valli and the symbolic languages of the Hermetists, among the symbolism of the doctrine of Love and of similar and derived movements; comparisons that indicate a clue and perhaps a proof of the existence and continuity of an initiatic tradition that arose in the Middle Ages. Unlike Valli, we however have several reservations about the purity of the Christian character of such tradition.

When one begins to recognize the existence of a ``false appearance" in a secret organization, beginning by degrees, it is right to doubt that if love and a noble [gentile] heart are one thing, the word ``gentile" can also have the meaning that it has in latin: sangue gentile [native blood]; and if Dante takes from Virgil the beautiful style, Virgil can also represent pagan initiation. But we will have the chance to return to these problems; for now, we limit ourselves to note how Boccaccio, who, Valli shows us, glorified the Templars, the same Boccaccio, author of a Genealogy of the Gods in the tenth story of the Decameron, makes jokes of the resurrection of the flesh, typical, i.e., of that same teaching that the Athenians mocked, saying to Saint Paul: ``we will hear you about this another time". We recall, in regards to Boccaccio, that in his third story he has Melchizedek say that between Judaism, Christianity, and Islam, ``no one knows which is the true faith". That Boccaccio puts phrases of this type even in Melchizedek's mouth, who occupies a position of the first order in tradition and in the esoteric hierarchy, is something that can make us reflect and can make us suspect what was the only Phoenix that with Zion joined the Appenines, as a sonnet says that goes under the name of Cino da Pistoia.

One final observation: in our earlier writing on the Knowledge of the Symbol, we had the chance to cite a passage from the Convivio, in which it shows how, according to Dante, the meanings to consider in the allegorical language were four, corresponding perhaps to the four degrees of the rite and of the organization. Of these four meanings, the most important for us, is the last, i.e., the anagogical meaning. Naturally this spiritual meaning, which is related to the process of spiritual development, cannot be understood and sometimes simply imagined, without the personal experience of it: who does not experience it, cannot understand it, says Dante. And it is for this reason that it almost always eludes those who up until now were occupied with secret language of the Fedeli d'Amore, unlike the meaning that we will call synagogic.

For example, to sleep means allegorically to live in ignorance, in the inertia of the intellect; morally, it means not to participate in the work of the organization; anagogically, it is the state opposite to that of initiatic Awakening. Valli thinks that, while the Vita Nuova was written in code, Dante abandoned in the Comedy the sectarian jargon; but if this is true, in part at least, through its moral or political meaning, since in the sacred poem the hostility against the Church is explicit and even extreme, it is not true for the anagogic meaning. This meaning is still and necessarily hidden under the veil of symbolism, in order to interpret it, it is necessary to possess the experience of the stages of consciousness to which it refers, and the knowledge of the symbols traditionally adopted to indicate them. For this reason, the true and higher meaning of the secret language of Dante and of the Fedeli d'Amore remains and will always remain a mystery to all those who ``sleep" and will continue to sleep.


\hfill



\flrightit{Posted on 2013-04-06 by Aeneas }

\begin{center}* * *\end{center}

\begin{footnotesize}\begin{sffamily}



\texttt{David on 2013-04-06 at 22:22 said: }

Few things. First, when they speak of potential intellect, is this in the same vein as Eckhart potentiality ? Second, what is the link between Dante and the Templars ? Thirdly, everything about this is more then welcome; it is very interesting and I would like to here more between Dante and Virgil, Fedeli d'Amore and esoterism. Thank you for the translation.


\hfill

\texttt{Cologero on 2013-04-06 at 22:45 said: }

The active intellect is a teaching from Aristotle. There was a lively debate leading up to Thomas on this topic, particularly from the Islamist philosophers Averroes and Avicenna. It is a large topic and the Catholic Encyclopedia, the Stanford Encyclopedia of Philosophy, and Wikipedia are some resources. Of course, we would be more interested in the esoteric meaning, which is why I mentioned Tomberg's second letter. That can be read as the active and passive intellects. The next question is why that is described as ``Love". This topic will come up again when I review a book on the relationship between the ideas of Thomas Aquinas and Dante, something which seems to have eluded Reghini.


\hfill

\texttt{Jason-Adam on 2013-04-08 at 14:54 said: }

It is possible to interpret this differently without resorting to paganism – one can see Dante as Dugin does as campaigning for the restoration of the Roman-Byzantine symphony of powers between the Pope and the Emperor that the harlot (the post-schism papacy) broke by trying to usurp the imperial rights. I do not share this view but am mentioning it just to show the variety of possibilities.


\hfill

\texttt{Cologero on 2013-04-08 at 20:17 said: }

Reghini's attempt to ``prove" its paganism is off base. Aside from the obvious fact that no pagans are rushing to incorporate Dante, there is the other uncomfortable (for Reghini) fact that much of Dante's Divine Comedy is related to Persian and other Islamic sources. The name ``alchemy" itself is of Arabic origin. Without going into all the details at this time, Rene Guenon's two essays on Valli's book, included in ``Insights into Christian Esoterism" should be consulted alongside Reghini's essay. Rather than marking a sharp departure from the spiritual tradition of the Middle Ages, Guenon's conclusion is more rational:

\begin{quotex}
\textbf{Dante's work, far from being contrary to the spirit of the Middle Ages, is one of its most perfect syntheses}, in the same rank as that of the cathedral builders; and the simplest initiatic facts enable us to understand without difficulty that there are very profound reasons for this correspondence. 

\end{quotex}

\hfill

\texttt{Cologero on 2013-04-18 at 23:02 said: }

\begin{quotex}
During the journey, which lasted thirteen days, the Magi took neither rest nor food; they did not feel any need, and this period seemed to them not to last more than a day. The nearer they came to Bethlehem, the brighter the star shone. It had the form of an eagle, flying through the air and moving its wings. Above it was a cross.

\end{quotex}
(From the Apocrypha, book ii.)


\hfill

\texttt{argusandphoenix on 2019-04-08 at 21:20 said: }

Dante just grows and grows as an iconic figure. Thank you, sir.


\hfill

\texttt{Cologero on 2020-04-06 at 07:44 said: }

Saint Bernard is the link between Dante and the Templars.


\hfill

\texttt{Graham on 2013-04-09 at 18:32 said: }

It seems like it is advisable to read the Comedy and the Convivio together, to understand the symbolism of Beatrice – perhaps with Guenon's volume near to hand.

Then again, if Reghini says one needs experience of the higher states of consciousness in order to comprehend the anagogic meanings, then all this reading could be futile.

How does one approach an initiatic text? Is there a method for reading a text in its depth? Something akin to lectio divina?


\hfill

\texttt{David on 2013-04-11 at 20:52 said: }

I am nowhere near being able to answer your question, or even understand in depth the Divine Comedy; but I think we cannot forget the sheer beauty and motivation of such a text (I'm not talking about pure emotionnal esthetics or poetry, but rather the drive to attain those states). I think those texts were written in part for this; if not, what reason is there to describe something only other people who attainted them would be able to understand ? It's like telling someone his own name. To me those text were both written as a fact and a motivation; but mostly as the bridge between both of them.


\hfill

\texttt{Cologero on 2020-08-14 at 07:24 said: }

Dante relied on the following texts that describe the path to higher states of consciousness.

Richard of Saint Victor, De Contemplatione

Saint Bernard, De Consideratione

Saint Augustine, De Quantitate Animae

For a review of De Quantitate Animae, see The Greatness of the Soul\footnote{\url{https://www.gornahoor.net/?p=12616}}


\end{sffamily}\end{footnotesize}
